% =====================================================================
% Appendix C -- Scope Evolution Timeline
% =====================================================================
\documentclass[../preamble.tex]{subfiles}
\begin{document}

\chapter{Scope Evolution Timeline}
\label{app:scope-evolution}

The thesis reports a narrower experimental scope than the original FYDP
proposal. The deviation is recorded explicitly here so that future readers
can interpret the methodology and results chapters accurately.

\section{Proposal-to-execution timeline}
\label{app:scope-timeline}

\begin{table}[htbp]
  \centering
  \small
  \caption{Evolution from the original FYDP proposal to the executed pipeline.
           ``Executed reality'' refers to the artefacts under
           \texttt{results/} and the figures produced by the scripts under
           \texttt{scripts/}.}
  \label{tab:scope-evolution}
  \begin{tabularx}{\linewidth}{@{}lXX@{}}
    \toprule
    \textbf{Item} & \textbf{Proposal claim} & \textbf{Executed reality} \\
    \midrule
    Preprocessing grid & nine backbones $\times$ six conditions
                       & nine backbones $\times$ two conditions (SRAD vs.\ Gaussian), with CLAHE fixed \\
    Datasets          & Figshare only
                       & Figshare + external PCOSgen (Sundari et~al., 2025) \\
    XAI methods       & Grad-CAM, LRP, SHAP
                       & Grad-CAM only \\
    Uncertainty       & MC-Dropout + uncertainty-gated referral
                       & MC-Dropout entropy + per-coverage statistics; referral remains a design concept \\
    Cross-validation  & 10-fold patient-level CV
                       & Single stratified 80/10/10 split on Figshare, single 2{,}560/640 split on PCOSgen train pool, separate held-out test \\
    Calibration       & ECE + temperature scaling
                       & Two-pass temperature scaling per-model and on the ensemble probability \\
    Architecture roster & ResNet, DenseNet, EfficientNet, InceptionV3 (4 families)
                       & ResNet, DenseNet, EfficientNet, ConvNeXt, MobileNetV3, ViT, Swin (7 families) \\
    Ensemble          & K-fold top-k finalist ensemble
                       & Pre-HPO probability-averaged 3-model ensemble (top-3 fine-tuned checkpoints) \\
    \bottomrule
  \end{tabularx}
\end{table}

\section{Why the scope changed}
\label{app:scope-why}

\subsection*{Preprocessing grid}

A nine-by-six grid would have required 54 foundation-stage runs, which would
have exceeded the available compute budget for a single project cycle.
Reducing the grid to two conditions isolates the denoiser as the variable
of interest and produces a directly interpretable comparison.

\subsection*{External dataset}

The original proposal stated Figshare only. PCOSgen became publicly available
in 2025 and its release was treated as an opportunity to evaluate
cross-dataset generalisation, which is a stronger test than a single
in-domain split. The added value of an external cohort outweighed the cost
of re-routing the fine-tune stage.

\subsection*{XAI methods}

The original brief listed three XAI methods. Running all three on every
fine-tuned checkpoint would have required substantial engineering effort for
each method, especially LRP under timm backbones. Grad-CAM was prioritised
because it integrates with the existing forward pass and because it was
sufficient to expose the padding shortcut. LRP and SHAP remain in the
designed-but-not-executed appendix.

\subsection*{Uncertainty-gated referral}

The brief proposed a deployed referral mechanism with a fixed entropy
cutoff. This thesis implements the per-coverage statistics from which a
deployment could set the cutoff, but does not deploy the mechanism. A
clinical deployment would require a defined referral cost, a target
coverage, and prospective validation, none of which are within scope.

\subsection*{Cross-validation}

A 10-fold cross-validation over the combined Figshare and PCOSgen corpus
would require a patient-level grouping that PCOSgen does not document. The
thesis therefore uses a single held-out test set with bootstrap confidence
intervals, and recommends nested cross-validation as future work.

\subsection*{Architecture roster}

InceptionV3 was removed because the contemporary transformer (ViT) and
mobile (MobileNetV3) families added complementary inductive biases that
InceptionV3 did not. Swin-Tiny was added to balance the parameter-count
envelope against ViT-Base/16.

\subsection*{Ensemble construction}

A $k$-fold top-k finalist ensemble is a more principled construction, but
the small PCOSgen validation set limits how much $k$-fold gains. The
pre-HPO probability-averaged ensemble is reported instead, with the
post-HPO retrain comparison included to show why the pre-HPO choice is
preferred.

\end{document}